% ==============================================================================
% PAKET-PAKET LATEX (PACKAGES)
% ==============================================================================

% Pengaturan Margin & Ukuran Kertas
\usepackage[
    a4paper,
    left=4cm,
    top=3cm,
    right=3cm,
    bottom=3cm,
    footskip=1.5cm,
    headheight=15pt
]{geometry}

% Tipografi: Times New Roman 12pt & Justifikasi Teks Rapi
\usepackage{mathptmx}
\usepackage[utf8]{inputenc}
\usepackage[T1]{fontenc}
\usepackage{microtype}

% Spasi dan Indentasi
\usepackage{setspace}
\usepackage{indentfirst} % Paragraf pertama setelah heading tetap menjorok masuk
\usepackage{enumitem}    % Pengendalian presisi spasi daftar (list/enumerate/itemize) tanpa space after/before

% Simbol & Rumus Matematika
\usepackage{amsmath}
\usepackage{amssymb}
\usepackage{amsfonts}

% Gambar dan Grafik
\usepackage{graphicx}
\usepackage{float}
\usepackage{tikz}
\usetikzlibrary{shapes.geometric, arrows.meta, positioning, fit, calc}

% Tabel Akademik
\usepackage{booktabs}
\usepackage{tabularx}
\usepackage{multirow}
\usepackage{array}
\usepackage{colortbl}

% Format Judul Bab & Subbab
\usepackage{titlesec}
\usepackage{etoolbox}

% Header dan Footer (Penomoran Halaman)
\usepackage{fancyhdr}

% Keterangan Gambar & Tabel
\usepackage{caption}
\usepackage{subcaption}

% Penyorotan Kode Program (Syntax Highlighting)
\usepackage{listings}
\usepackage{xcolor}

% Pengaturan Daftar Isi, Gambar, dan Tabel
\usepackage{titletoc}

% Sitasi & Daftar Pustaka Gaya APA (American Psychological Association)
\usepackage[natbibapa]{apacite}

% Hyperlink dokumen PDF
\usepackage{hyperref}
\usepackage{url}
